% Shared preamble of the papers in docs/paper (paper.tex, jlt/jlt.tex).
% Each paper adds its own \addbibresource (the path to references.bib is
% relative to the paper's directory).
% LuaLaTeX handles UTF-8 natively; fontspec lets the document pick up any
% Unicode glyphs present in the default font, including mathematical script
% characters used in listings (ℕ, ℤ, ℚ, ℝ, ℂ, Ø, and similar). We keep the
% `newunicodechar` fallbacks for body text so amsmath glyphs (\mathbb) are
% used in paragraphs.  Code listings use the `minted` package (Pygments-
% backed) with native UTF-8 — no per-character literate table needed.
\usepackage{fontspec}
\usepackage{newunicodechar}
\usepackage[english]{babel}    % Language support
\usepackage{geometry}          % Adjust page margins
\geometry{margin=1in}

% --- Recommended Math & Science Packages ---
\usepackage{amsmath, amssymb, amsthm} % Essential math tools
\usepackage{mathtools}
\usepackage{mathrsfs}

% Theorem environments (amsthm is loaded above).
\newtheorem{theorem}{Theorem}
\newtheorem{definition}{Definition}
\newtheorem{corollary}{Corollary}
\newtheorem{lemma}{Lemma}
\newtheorem{proposition}{Proposition}
\newtheorem{conjecture}{Conjecture}
\usepackage{graphicx}                  % Support for images
\usepackage{hyperref}                  % Clickable links & cross-refs
\usepackage{booktabs}                  % High-quality tables
% csquotes is loaded later — after `minted` (which pulls in `fvextra`)
% to silence the fvextra/csquotes load-order warning.

\newunicodechar{∩}{\ensuremath{\cap}}
\newunicodechar{ℚ}{\ensuremath{\mathbb{Q}}}
\newunicodechar{ℝ}{\ensuremath{\mathbb{R}}}
\newunicodechar{𝔽}{\ensuremath{\mathbb{F}}}
\newunicodechar{∃}{\ensuremath{\exists}}
\newunicodechar{π}{\ensuremath{\pi}}
\newunicodechar{⊂}{\ensuremath{\subset}}
\newunicodechar{⊆}{\ensuremath{\subseteq}}
\newunicodechar{∈}{\ensuremath{\in}}
\newunicodechar{∉}{\ensuremath{\notin}}
\newunicodechar{∧}{\ensuremath{\land}}
\newunicodechar{∨}{\ensuremath{\lor}}
\newunicodechar{≤}{\ensuremath{\leq}}
\newunicodechar{≥}{\ensuremath{\geq}}
\newunicodechar{μ}{\ensuremath{\mu}}
\newunicodechar{α}{\ensuremath{\alpha}}
\newunicodechar{β}{\ensuremath{\beta}}
\newunicodechar{λ}{\ensuremath{\lambda}}
\newunicodechar{Ω}{\ensuremath{\Omega}}
\newunicodechar{𝔸}{\ensuremath{\mathbb{A}}}
\newunicodechar{↔}{\ensuremath{\leftrightarrow}}
\newunicodechar{χ}{\ensuremath{\chi}}
\newunicodechar{Σ}{\ensuremath{\Sigma}}
\newunicodechar{Τ}{\ensuremath{\mathrm{T}}}

% For code listings
\usepackage{xcolor}
% Landscape tables:
\usepackage[figuresright]{rotating}
\usepackage[table]{xcolor}
\usepackage{pdflscape}
\usepackage{tabularx}
\usepackage{booktabs}


% Graphs
\usepackage{graphviz}

\usepackage{tikz}
\usetikzlibrary{shapes.geometric, arrows.meta, positioning, calc, backgrounds}
\usepackage{tikz-cd}

\usepackage{draftwatermark}
\SetWatermarkScale{5}

\usepackage[style=numeric, sorting=none]{biblatex}

% `minted` (Pygments-backed) handles UTF-8 natively in code listings —
% no per-character literate table needed.  Requires `lualatex -shell-escape`
% (the build pipeline in docs/paper/Makefile and the CI install step in
% the project Makefile are updated accordingly).
\usepackage{minted}
\usepackage{csquotes}  % Load after minted/fvextra to silence load-order warning.

% Use JuliaMono for code listings — designed specifically for Unicode-rich
% source (Julia uses the full math repertoire) with native coverage of
% Greek (φ, Σ, Τ, χ, ...), math double-struck (ℕ, ℤ, ℚ, ℝ, ℂ, 𝔻, 𝔽, ...),
% math-fraktur, math-script, sub/superscripts, and arrows.  Replaces the
% earlier `lmmono10` (no Greek) and `Fira Mono` (no math-double-struck)
% choices made in the lstlisting era.  Installed locally via
% `brew install --cask font-juliamono`; CI install adds `fonts-juliamono`
% (Ubuntu 22.04+) under ci-install-paper-deps in the project Makefile.
% Scale=0.90 keeps glyph metrics close to the document's serif body.
\setmonofont{JuliaMono}[Scale=0.90]

\setminted{
  fontsize=\small,
  breaklines=true,
  breakanywhere=true,
  frame=single,
  framesep=2mm,
  autogobble,
}
\usemintedstyle{friendly}

% C++ is the default code-listing language in this paper; declare short
% block + inline aliases so source markup reads `\begin{cpp}...\end{cpp}`
% and `\cppinline{...}` rather than the verbose `\begin{minted}{cpp}` /
% `\mintinline{cpp}{...}` forms.  The optional `[name]` first argument
% names the environment directly; without it minted would produce
% `\begin{cppcode}` and `\cpp{...}`, which would collide with our \cpp
% short alias usage.
\newminted[cpp]{cpp}{}
\newmintinline[cppinline]{cpp}{}

\definecolor{codegreen}{rgb}{0,0.6,0}
\definecolor{codegray}{rgb}{0.5,0.5,0.5}
\definecolor{codepurple}{rgb}{0.58,0,0.82}

\newcommand{\bind}{\mathbin{>\!\!>\mkern-6.7mu=}}
\newcommand{\fish}{\mathbin{\gg}}
\newcommand{\Cpp}[1]{C\kern-.05em\raise.3ex\hbox{\footnotesize ++}#1}
